
\documentclass[12pt,a4paper]{article}

\usepackage[T1]{fontenc}
\usepackage[utf8]{inputenc}
\usepackage[polish]{babel}
\usepackage{lmodern}
\usepackage{geometry}
\usepackage{amsmath}
\usepackage{graphicx}

\geometry{margin=2.5cm}

\title{Sprawozdanie z analizy danych}
\author{Jakub Muszka}
\date{\today}

\begin{document}

\maketitle

\tableofcontents
\newpage

\section{Liczności gatunków irysów}

\begin{table}[htbp]
    \centering
    \caption{Liczebność poszczególnych gatunków irysów}
    \label{tab:wyniki}
    \begin{tabular}{|c|c|}
        \hline
        Gatunek & Liczebność(\%) \\
        \hline
        Setosa & 50(33.33\%) \\
        Versicolor & 50(33.33\%) \\
        Virginica & 50(33.33\%) \\
        Razem & 150(100\%) \\
        \hline
    \end{tabular}
\end{table}

\section{Miary rozkładu cech irysów}

\begin{table}
    \caption{Miary rozkładu dla cech irysów}
    \label{tab:miary}
    \resizebox{\textwidth}{!}{
    \begin{tabular}{|c|c|c|c|c|c|}
        \hline
        Cecha & Minimum & Śr. arytm. (± odch. stand.) & Mediana (Q1 - Q3) & Maksimum \\
        \hline
        Długość działki kielicha & 4.3 & 5.84 ± 0.83 & 5.8 ((5.1, 6.4)) & 7.9 \\
        Szerokość działki kielicha & 2 & 3.06 ± 0.44  & 3.0 ((2.8, 3.3)) & 4.4 \\
        Długość płatka & 1 & 3.76 ± 1.77 & 4.35 ((1.6, 5.1)) & 6.9 \\
        Szerokość płatka & 0.1 & 1.20 ± 0.76  & 1.3 ((0.3, 1.8))  & 2.5 \\
        \hline
    \end{tabular}
    }
\end{table}

\section{Charakterystyka rozkładu cech irysów}

\section{Współzależności pomiędzy cechami irysów}




\end{document}